\section{Conclusão}

Este artigo apresentou uma versão reestruturada e mais explicitamente científica do Sonar Municipal, um pipeline para descoberta de instâncias do SAPL, extração de PLs, tradução de ementas em ações, busca semântica e associação exploratória com indicadores oficiais municipais. Na execução congelada de 12/11/2025, o sistema encontrou 1.259 instâncias do SAPL, extraiu PLs de 322 delas e reuniu um acervo de 220.065 proposições. O tradutor \textit{ementa} $\rightarrow$ \textit{ação} obteve BERTScore de 84\% em validação, e a transformação do acervo completo em ações mostrou viabilidade operacional.

A análise exploratória do dataset mostrou que a base final é suficientemente rica para fortalecer o artigo com evidências adicionais de cobertura, qualidade e compressão textual: o acervo cobre 19 UFs, concentra 43,63\% dos registros entre 2021 e 2025 e reduz o comprimento médio das ementas para cerca de 56\% a 59\% após a textualização em ações. Além disso, o estudo de caso de violência contra a mulher mostrou que o pipeline já consegue sair do nível abstrato e produzir consultas, recomendações e clusters interpretáveis em escala intermunicipal.

Ao mesmo tempo, a reorganização do artigo deixa claro que a principal fronteira de amadurecimento não está mais na existência do pipeline, mas na robustez da sua validação. Nesta rodada, o texto assume explicitamente a ausência de benchmark formal para recuperação e a inexistência de um conjunto de teste independente para o tradutor. Os próximos passos mais imediatos, portanto, são menos arquiteturais e mais editoriais e experimentais: consolidar a bibliografia final, decidir se haverá \textit{baselines} e avaliação humana adicional para o tradutor e incorporar uma captura da interface principal na versão pronta para submissão.
